\chapter{补充内容}

附录是与论文内容密切相关、但编入正文又影响整篇论文编排的条理和逻辑性的资料，例如某些重要的数据表格、计算程序、统计表等，是论文主体的补充内容，可根据需要设置。

\begin{figure}[htbp]
  \centering
  \includegraphics[width=0.42\linewidth]{./content/test.png}
  \caption{fig-appendix}
  \label{fig:fig-appendix}
\end{figure}

以下为分图示例以下为分图示例

以下为分图示例以下为分图示例

以下为分图示例以下为分图示例

以下为分图示例以下为分图示例以下为分图示例以下为分图示例

以下为分图示例：

\begin{figure}[htbp]
  \centering
  \subcaptionbox{分图 A\label{fig:subfig-a}}
    {\includegraphics[width=0.4\linewidth]{./content/test.png}}
  \subcaptionbox{分图 B\label{fig:subfig-b}}
    {\includegraphics[width=0.4\linewidth]{./content/test.png}}
  \caption{多个分图的示例}
  \label{fig:multi-image}
\end{figure}

以下为表格示例

以下为表格示例

以下为表格示例

以下为表格示例

以下为表格示例

以下为表格示例以下为表格示例

以下为表格示例：

\begin{table}[htbp]
  \centering
  \caption{table-example-appendix}
  \begin{tabular}{cccccc}
    \toprule
    数据库  & 原始图像数 & 失真图像数 & 失真类型 & 主观评价法 & 参评人员数 \\
    \midrule
    LIVE    & 29         & 779        & 5       & DMOS       & 161        \\
    TID2013 & 25         & 3000       & 24      & MOS        & 971        \\
    \bottomrule
  \end{tabular}
  \label{tab:table-example-appendix}
\end{table}

以下为代码示例
以下为代码示例

以下为代码示例

以下为代码示例

以下为代码示例以下为代码示例

以下为代码示例：

\begin{listing}[htbp]
  \begin{minted}[baselinestretch=1.0,linenos=true,breaklines,autogobble]{c}
    /**
     * @brief  LED亮度控制中断
     * @param  none
     * @retval none
     */
    typedef uint8_t  LED8BitsType;
    void LED_PeriodElapsedCallback(void){
        int i = 0;
        static LED8BitsType led_bright_tick = 0;
        led_bright_tick += 1;
        if(led_bright_tick>=100){
            led_bright_tick = 0;
        }
        if(led_bright_tick == led.brightness){
            HAL_GPIO_WritePin(FnLED_GPIO_Port, FnLED_Pin, GPIO_PIN_RESET);
        }else if(led_bright_tick == 0){
            HAL_GPIO_WritePin(FnLED_GPIO_Port, FnLED_Pin, GPIO_PIN_RESET);
        }
    }
    \end{minted}
    \caption{code-example}
    \label{listing:example}
\end{listing}

